\documentclass[10pt,a4paper]{amsart}
\usepackage[margin=19mm]{geometry}
\usepackage{amsmath,amssymb,mathtools}
\usepackage[scr=rsfs,cal=euler]{mathalfa}
\usepackage{xfrac}
\usepackage[unicode,hidelinks,bookmarks=false]{hyperref}
\newcommand{\ie}{i.e.\ }
\newcommand{\cf}{cf.\ }
\newcommand{\resp}{respectively\ }
\newcommand{\Subsection}{\subsection}
\providecommand{\nobreakdash}{\nobreak\hbox{-}}
\setlength{\emergencystretch}{2em}
\multlinegap0pt
\usepackage{bm}
\usepackage{bbm}
\usepackage{dsfont}
\usepackage{xspace}
\usepackage{cancel}
\usepackage{mathtools}
\usepackage[shortlabels]{enumitem}
\usepackage{amsthm}
\makeatletter
\@ifundefined{lemma}{\newtheorem{lemma}{Lemma}}{}
\@ifundefined{proposition}{\newtheorem{proposition}[lemma]{Proposition}}{}
\makeatother
\DeclareMathOperator{\Eff}{\overline{Eff}}
\DeclareMathOperator{\Mov}{Mov}
\DeclareMathOperator{\Nef}{Nef}
\newcommand{\bb}{\ensuremath{\mathfrak{b}}}
\renewcommand{\geq}{\geqslant}
\title[Fujita-Zariski decompositions on some product threefolds]{Fujita-Zariski decompositions on some product threefolds}
\author{Mihai Fulger, Victor Lozovanu}
\date{Source: arXiv:2506.17956v1 (CC BY 4.0)}
\begin{document}
\maketitle

\noindent\emph{Source and licence.} Fujita-Zariski decompositions on some product threefolds. Version arXiv:2506.17956v1 (CC BY 4.0), distributed under the Creative Commons Attribution 4.0 International licence, as recorded in the arXiv licence metadata for this exact version and shown on the abstract page https://arxiv.org/abs/2506.17956v1. Full rights notice: licenses/arxiv-2506.17956-CC-BY-4.0.txt.\par\smallskip

\noindent\textit{Editorial scope.} Counted unit: the Proposition whose source label is \texttt{prop:blCxP2positive} in the article (source0.tex lines 543--557, source label \texttt{prop:blCxP2positive}), delivered together with the article's own complete original proof of it (source0.tex lines 558--578, one \texttt{proof} environment of 21 lines). No mathematics is written, repaired or re-derived: statement and proof are the article's own text line for line. Required context retained uncounted: none. Every definition, display and label of those lines is retained unchanged. Omitted from this package and not redistributed: 1--542, 579--1837. Nothing inside a retained excerpt is omitted or altered: every retained body is delivered exactly as the article's lines are written, so no omission receipt of any kind accompanies this package. Citations: the retained excerpts carry no citation command at all, so no bibliography entry of the article is transcribed and no reference is redistributed. Reference adaptation: none was needed and none was made; every reference inside the delivered unit resolves inside it, so no unresolved reference or citation is produced. Printed slips kept literal: no printed word, symbol, hypothesis or number of the counted statement or of its proof was altered. Layout and class: the article's own class is \texttt{amsart} with packages including latexsym, amsmath, amssymb, amsthm, mathrsfs, wasysym, mathtools, paralist, caption, subcaption, fontenc, lmodern, hyperref, amsrefs; the wrapper is the lane's pinned \texttt{amsart} editorial wrapper at 10pt on A4, which restates the theorem-like environments the retained text needs and adds the article's own macro definitions that the retained text needs (\texttt{\textbackslash Eff}, \texttt{\textbackslash Mov}, \texttt{\textbackslash Nef}, \texttt{\textbackslash bb}, \texttt{\textbackslash geq}), with extra wrapper packages (bm, bbm, dsfont, xspace, cancel, mathtools, enumitem). The delivered numbering is the unit's own. Assets: no retained span contains a figure environment, a graphics-inclusion command, a PDF-page inclusion, a table or a TikZ picture; no image file is copied into this package, the package declares no asset dependency and no image file is redistributed with it. Licence: version arXiv:2506.17956v1 of this article is distributed under the Creative Commons Attribution 4.0 International licence, as recorded in the arXiv licence metadata for this exact version and shown on the abstract page https://arxiv.org/abs/2506.17956v1. A full rights notice is delivered at licenses/arxiv-2506.17956-CC-BY-4.0.txt. Characters: every retained excerpt is pure ASCII, so the delivered engine, XeLaTeX with its default Latin Modern text font, reports no missing character.
\begin{proposition}\label{prop:blCxP2positive}
	With notation as above,
	\begin{enumerate}[(1)]
		\item $\Eff(\overline X)=\langle\overline f,\ \overline H,\ E\rangle$. In particular $\mu(L;x)=a+b$.
		\smallskip
		
		\item $\Mov(\overline X)=\langle \pi^*f,\ \pi^*H,\ \overline H\rangle$.  In particular $\nu(L;x)=b$.
		\smallskip
		
		\item $\Nef(\overline X)=\langle \pi^*f,\ \pi^*H,\ \overline f+\overline H+E\rangle$. In particular $\epsilon(L;x)=\min\{a,b\}$. 
		\smallskip
		
		\item For $t\in[0,a+b)$, we have $P_{\sigma}(L_t)=L_t-(t-b)_+\cdot\overline f=\min\{a,a+b-t\}\cdot\overline f+b\overline H+(a+b-t)E$.
	\end{enumerate}
\end{proposition}

\begin{proof} 
Clearly $N^1(\overline X)_{\bb R}$ is generated by $\pi^*f$, $\pi^*H$, and $E$.

(1). Let $\alpha=a\overline f+b\overline H+cE$ be the class of an irreducible effective divisor on $\overline X$ distinct from $\overline f$, $\overline H$, and $E$. It is sufficient to prove that $a,b,c\geq 0$. Its pushforward to $X$ is also an irreducible effective divisor, which gives $a,b\geq 0$.
Its restriction to $\overline f\simeq{\rm Bl}_p\bb P^2$ has class $bh+(c-a-b)E$ where $h$ is the class of the pullback of a line in $\bb P^2$ and by abuse $E$ is the exceptional line over $p$. This is effective by our assumptions, which implies $c\geq a\geq 0$.

(2). The pullback of a nef divisor class is again nef, and in particular movable. Furthermore $\overline H\equiv\overline H'$. This implies that $\overline H$ is movable. The claimed generators are indeed movable. Conversely, let $\alpha=a\overline f+b\overline H+cE$ be a movable divisor class. It is in particular pseudoeffective, thus $a,b,c\geq 0$. Its restriction to any divisor must also be pseudoeffective. 
The restriction to $E=\bb P^2$ has degree $a+b-c$, thus $a+b\geq c$.
The restriction to $\overline f$, as in $(1)$ gives the condition $c\geq a$.
The cone cut out by the conditions $a+b\geq c\geq a\geq 0$ is the cone in the conclusion. 

(3). The dual of the cone of curves $\langle\overline C_x,\ \overline{\bb P}^1_x,\ \ell\rangle$ is generated by  $\pi^*f$, $\pi^*H$, and $\overline f+\overline H+E$ as one checks directly. It is then sufficient to prove that these three claimed generators are nef. This is clear for the first two. For the last one, which can be written as $\overline f+\pi^*H$, it is sufficient to prove that its restriction to $\overline f$ is nef. This restriction is $h-E$, which is indeed nef. 

Dually, the Mori cone of curves is generated by $\overline C_x$, $\overline{\bb P}^1_x$, and $\ell$. Thus a class $\alpha$ as above is nef if and only if $a+b\geq c\geq\max\{a,b\}$.


(4). Negative parts (e.g., $N_{\sigma}(L_t)$) cannot contain movable components like $\overline H$. Since $L$ is ample, then $x\not\in{\bf B}_+(L)$ and then $E$ is not a component of any ${\bf B}_+(L_t)$ while $L_t$ is big, in particular it is also not in the support of $N_{\sigma}(L_t)$. We deduce that $N_{\sigma}(L_t)=m_t\cdot\overline f$ for some $m_t\geq 0$. 
Then $(P_{\sigma}(L_t))|_{\overline f}=(L_t-m_t\cdot\overline f)|_{\overline f}\equiv b(h-E)+(b-t+m_t)E$ is pseudoeffective which implies $b-t+m_t\geq 0$, i.e., $m_t\geq(t-b)_+$.

It follows easily from (2) that the claimed formula for the positive part is indeed movable. The maximality of the positive part $P_{\sigma}(L_t)$ among movable divisors dominated by $L_t$ implies the reverse inequality $a-m_t\geq\min\{a,a+b-t\}$.  
\end{proof} 

\end{document}
